\documentclass[instructions.tex]{subfiles}
\begin{document}
 
\chapter{Pour être pénitent, il faut vivre dans l'humilité.}
 
Le Sage nous apprend que l’orgueil est le principe de tout péché\footnote{Initium omnis peccati superbia. Sap.} ; parce que tout pécheur s’élève contre Dieu, et s’éloigne de lui, il faut donc que l’humilité le rapproche de Dieu. De quelque côté que le pécheur se considere, soit du côté de Dieu, soit du côté de lui-même, soit du côté des autres, tout doit l’humilier.
 
\primo Du côté de Dieu, il doit se regarder comme un esclave révolté, comme un criminel de lèse-majesté, comme un perfide et un ingrat qui a trahi son maître; comme un enfant rebelle et un orgueilleux coupable, qui a déshonoré son père et son Dieu. Oui, vous avez déshonoré Dieu, et vous l’avez plus déshonoré par un seul péché mortel, que jamais vous ne pourriez l’honorer, quand vous feriez autant de bonnes œuvres que tous les justes. Quel sujet de confusion! et comment oseriez-vous vous prévaloir de vos bonnes œuvres?
 
\secundo Du côté de vous-même. O homme! qu’êtes-vous devenu par le péché? Vous avez donné le coup de la mort à votre âme, vous l’avez défigurée; vous lui avez imprimé, selon l’Ecriture, le caractère de la bête, parce que vous n’avez suivi que vos sens et vos penchans, comme les bêtes\footnote{Comparatus est jumentis insipientibus, Ps.}. O mon Dieu! disait David, comment oserai-je paraître devant vous, après m’être dégradé jusqu’à devenir à vos yeux comme une bête de charge.
 
Ce fut un état bien humiliant pour le Sauveur, de vivre dans le désert parmi les bêtes farouches. Mais reconnoissez avec humilité que vous ne méritez pas même d’être en présence de Jésus-Christ dans le rang des bêtes sauvages : elles n’ont jamais péché, elles ne sont capables ni de connoître, ni d’aimer Dieu; au lieu que vous, qui n’êtes fait que pour l’aimer, vous ne vous servez de votre raison et de ses bienfaits que pour l’offenser.
 
Vous devriez donc vous placer au-dessous des animaux, au-dessus même des démons. Il y en a un parmi vous qui est un démon, disait le Sauveur en parlant de Judas. Voilà ce que vous devriez penser de vous-même, avec cette différence, que le démon n’a péché qu’une fois; et vous, de combien de crimes multipliés n’êtes-vous pas coupable! Pouvez-vous, après cela, vous estimer? On n’estime point ce qui est défiguré et horrible : or, pouvez-vous voir quelque chose de plus horrible que votre conscience? Y a-t-il donc un objet que vous deviez plus mépriser que vous-même?
 
\tertio Du côté des autres. Vous oubliez que vous êtes pécheur, lorsque vous les méprisez. Eussiez-vous la vertu d’un saint, votre conscience doit vous faire dire, avec plus de raison que saint Paul : Je suis le plus grand des pécheurs, le plus misérable de tous les hommes; parce que vous devez avoir plus de confusion et d’horreur d’un seul de vos péchés, que des péchés de tous les autres.
 
D’ailleurs, de quoi n’êtes-vous pas capable? Un scélérat eût-il commis des crimes que vous ne commites jamais, croyez qu’il est encore meilleur que vous, parce qu’il n’a reçu ni tant de lumieres, ni tant de grâces, et que vous auriez vous-même commis des crimes plus énormes, si Dieu vous avait abandonné à vous-même. Le fond de notre cœur est si mauvais, que nous devons, en quelque sorte, regarder comme nôtres les péchés mêmes que nous n’avons pas faits; ce qui a fait dire à saint Augustin : Seigneur, vous m’avez fait autant de grâce en me préservant des péchés que je n’ai pas commis, qu’en me pardonnant ceux que j’ai faits\footnote{Omnia mihi dimissa fateor, et quæ meà sponte fecit mala, et quæ, te duce, non feci. Confess., lib. 10. cap. 10.}.
 
La vue de votre misere doit vous accompagner par-tout, vous reconnoissant indigne de la société des créatures, après avoir offensé leur Créateur. Quand vous êtes dans les assemblées de l’Eglise, souvenez-vous de cette parole de Job : Les enfans de Dieu s’étant assemblés, Satan se trouva avec eux. Soyez dans la confusion de vous trouver dans le lieu saint avec tant de fideles, dont plusieurs ont la pureté des anges, tandis que vous êtes un pécheur abominable; et croyez que c’est beaucoup pour vous, que l’Eglise vous souffre parmi ses enfans, vous qui devriez être sous les pieds des démons.
 
Si vous avez des richesses ou des emplois honorables, pensez que vous ne les avez pas mérités; que les autres en sont plus dignes que vous, et que votre élévation est peut-être un effet de la justice de Dieu sur vous. Reconnoissez qu’un pécheur comme vous, qui a mérité les anathêmes de toutes les créatures, ne mérite pas d’être élevé au-dessus des autres, ni de leur commander\footnote{Quandò magnus es, humilia te in omnibus. Eccli. 3. 20.}.
 
Si l’on vous méprise, si l’on vous persécute, que la vue de vos péchés vous fasse aussitôt souvenir que vous ne méritez ni l’estime ni l’attention de personne; qu’un pécheur qui mérite l’enfer ne peut être traité trop durement en cette vie; et qu’à bien prendre la chose selon Dieu, vous n’avez pas lieu de vous plaindre, dès qu’on vous traite comme vous le méritez.
 
Oh! si, à l’exemple du publicain, vous aviez ces sentimens d’humilité, quelque pécheur que vous ayez été, vous seriez un vrai pénitent et un grand saint, parce que Dieu ne rejette jamais un cœur contrit et humilié\footnote{Cor contritum et humiliatum, Deus, non despicies. Ps. 50. 19.}. Sans cet esprit d’humilité, on ne se sauve pas. Si vous ne devenez humble comme des enfans, dit Jésus-Christ, vous n’entrerez jamais dans le royaume des cieux.
 
\section{Résolutions}
 
\primo Je me considérerai désormais comme un néant révolté contre Dieu.—

\secundo Je me confondrai à mes propres yeux, considérant combien j'ai dégradé en moi l'ouvrage du Tout-Puissant.

\tertio Je me mettrai, dans ma propre pensée, au dessus des plus grands pécheurs.

\quarto Et je souffrirai les mépris comme une des moindres peines que j'ai méritées de la part des créatures. 

\prayer{ô mon Dieu ! n'écrasez pas ce ver de terre, qui a osé s'élever contre vous et contre votre loi sainte: Il reconnaît enfin sa bassesse et sa profonde misère.}
 
\end{document}
